% Part: incompleteness
% Chapter: incompleteness-provability
% Section: lob-thm

\documentclass[../../../include/open-logic-section]{subfiles}

\begin{document}

\olfileid{inc}{inp}{lob}

\olsection{لوب دا قضیہ}

نظریے~$\Th{T}$ دا گوڈل جملہ $\lnot
\OProv[\Th{T}](y)$ دا مقرر نقطہ اے،یعنی ایسا !!a{sentence}~$!G$
جیہڑا
\[
\Th{T} \Proves \lnot \OProv[\Th{T}](\gn{!G}) \liff !G.
\]
پورا کردا اے۔ جے نظریہ بے تضاد ہووے تاں ایہہ !!{derivable}
نہیں: جے $\Th{T} \Proves !G$ ہووے،تاں (a) !!{derivability}
شرط~(1) توں $\Th{T} \Proves \OProv[\Th{T}](\gn{!G})$،
تے (b) $\Th{T}
\Proves !G$ نال $\Th{T} \Proves \lnot \OProv[\Th{T}](\gn{!G})
\liff !G$ مل کے $\Th{T} \Proves \lnot
\OProv[\Th{T}](\gn{!G})$ دیندے نیں؛فیر $\Th{T}$
متضاد ہووے گا۔ ہن قدرتی سوال ایہہ اے کہ
$\OProv[\Th{T}](y)$ دے مقرر نقطے دی کی حیثیت اے،یعنی
ایسا !!a{sentence}~$!H$ جیہڑا
\[
\Th{T} \Proves \OProv[\Th{T}](\gn{!H}) \liff !H.
\]
پورا کردا اے۔ جے ایہہ !!{derivable} ہووے،تاں شرط~(1)
توں $\Th{T} \Proves \OProv[\Th{T}](\gn{!H})$؛پر اُتلے
ہم ارزی والے بیان اُتے قیاسِ استثنائی لائیے تاں وی ایہو نتیجہ
ملدا اے۔ ایس لئی گوڈل جملے والی دلیل وانگ اتھوں
$\Th{T}$ دا تضاد نہیں نکلدا،کم از کم اوسے دلیل نال تاں نہیں۔
البتہ ایس توں ایہہ وی ثابت نہیں ہوندا کہ $\Th{T}$
\emph{واقعی} $!H$ نوں !!{derive} کردا اے۔

اس سوال وچ پیش رفت لئی اس نوں کجھ عمومی روپ وچ ویکھدے آں۔
مقرر نقطے دی ہم ارزی دا کھبے توں سجے ول رخ،
$\OProv[\Th{T}](\gn{!H}) \lif !H$،اک عمومی خاکے
$\OProv[\Th{T}](\gn{!A}) \lif !A$ دی مثال اے جیہنوں
\emph{انعکاس دا اصول} کہندے نیں۔ ایہہ ناں ایس لئی اے کہ،
اک معنی وچ،ایہہ $\Th{T}$ دے اپنے !!{derive} کیتے ہوئیاں
جملیاں بارے غور نوں ظاہر کردا اے؛بنیادی گل ایہہ اے:
«جے $\Th{T}$، $!A$ نوں !!{derive} کر سکدا اے،تاں
$!A$ سچ اے»،ہر~$!A$ لئی۔ ایہ گل صرف اوہناں نظریاں لئی سچ اے جنہاں دے
ثابت کیتے سارے جملے سچے ہون۔ ایس توں اندازہ ہوندا اے کہ نظریے عام طور اُتے اس
اصول دی ہر مثال نوں !!{derive} نہیں کرن گے۔ تاں کِہڑیاں
مثالاں کوئی کافی قوی نظریہ،جیہڑا !!{derivability} شرطاں وی
پوری کردا ہووے،!!{derive} کر سکدا اے؟ یقیناً اوہ ساریاں
جنہاں وچ $!A$ خود !!{derivable} اے۔ اگلا نتیجہ دسدَا
اے کہ صرف ایہو مثالاں نیں۔

\begin{thm}\ollabel{thm:lob}
منّو $\Th{T}$، $\Th{Q}$ دی توسیع کرن والا کوئی
!!a{axiomatizable} نظریہ اے،تے
$\OProv[\Th{T}](y)$ اک ایسا فارمولا اے جیہڑا
\olref[2in]{sec} دیاں P1--P3 شرطاں پوری کردا اے۔ جے
$\Th{T}$، $\OProv[\Th{T}](\gn{!A}) \lif !A$ نوں
!!{derive}s کردا اے،تاں واقعی $\Th{T}$، $!A$ نوں
!!{derive}s کردا اے۔
\end{thm}

دوجے لفظاں وچ،جے $\Th{T} \Proves/ !A$،تاں
$\Th{T} \Proves/
\OProv[\Th{T}](\gn{!A}) \lif !A$۔ ایس نتیجے نوں
لوب دا قضیہ کہندے نیں۔

\begin{explain}
لوب دے قضیے دے ثبوت دا خیال سانتا کلاز دے موجود ہون دی
اک چترائی بھری دلیل توں ملدا اے۔ (جے ایہہ نتیجہ پسند نہ
آوے،تاں اپنی پسند دا کوئی ہور نتیجہ رکھ لو۔) دلیل ایہہ اے:
\begin{enumerate}
\item $X$ اوہ جملہ ہووے: «جے $X$ سچ اے،تاں سانتا کلاز
  موجود اے»۔
\item منّو $X$ سچ اے۔
\item فیر اس جملے دی کہی گل درست اے؛یعنی جے $X$ سچ اے
  تاں سانتا کلاز موجود اے۔
\item چونکہ اسیں $X$ نوں سچ منّیا اے،(2) تے~(3) توں
  قیاسِ استثنائی نال ایہہ نتیجہ کڈھ سکدے آں کہ سانتا کلاز
  موجود اے۔
\item اسیں فرض~(2)،«$X$ سچ اے»،توں (4)،«سانتا کلاز
  موجود اے»،اخذ کر لیا۔ شرطیہ ثبوت نال دکھایا: «جے $X$
  سچ اے،تاں سانتا کلاز موجود اے»۔
\item پر ایہہ تاں عین جملہ~$X$ اے۔ ایس لئی اسیں دکھا
  دتا کہ $X$ سچ اے۔
\item فیر اُتلی (2)--(4) والی دلیل نال سانتا کلاز موجود اے۔
\end{enumerate}
ایس خیال نوں رسمی بنا کے «سچ اے» دی تھاں «!!{derivable}
اے» تے «سانتا کلاز موجود اے» دی تھاں~$!A$ رکھیے،تاں
لوب دے قضیے دا ثبوت ملدا اے۔ چال ایہہ اے کہ مقرر نقطے دا
لمہ !!{formula}~$\OProv[\Th{T}](y) \lif !A$ اُتے
لائیے۔ ایس دا مقرر نقطہ پچھلے خاکے دے جملے~$X$ دے
مطابق اے۔
\end{explain}

\begin{proof}[\olref{thm:lob} دا ثبوت]
منّو $!A$ ایسا !!a{sentence} اے جیہڑے لئی $\Th{T}$،
$\OProv[\Th{T}](\gn{!A}) \lif !A$ نوں !!{derive}s
کردا اے۔ $!B(y)$ !!{formula}~$\OProv[\Th{T}](y)
\lif !A$ ہووے؛مقرر نقطے دے لمے نال ایسا !!a{sentence}~$!D$
لبھو جیہڑے لئی $\Th{T}$، $!D \liff !B(\gn{!D})$ نوں
!!{derive}s کردا اے۔ فیر ہِٹھلا ہر بیان
$\Th{T}$ وچ !!{derivable} اے:
\begin{align}
  & !D \liff (\OProv[\Th{T}](\gn{!D}) \lif !A) \ollabel{L-1}\\
  & \qquad \text{$!D$، $!B(y)$ دا مقرر نقطہ اے}\notag \\
  & !D \lif (\OProv[\Th{T}](\gn{!D}) \lif !A) \ollabel{L-2}\\
  & \qquad\text{\olref{L-1} توں}\notag\\
  & \OProv[\Th{T}](\gn{!D \lif (\OProv[\Th{T}](\gn{!D}) \lif !A)}) \ollabel{L-3}\\
  & \qquad \text{\olref{L-2} توں شرط P1 نال}\notag \\
  & \OProv[\Th{T}](\gn{!D}) \lif \OProv[\Th{T}](\gn{\OProv[\Th{T}](\gn{!D}) \lif !A})
  \ollabel{L-4}\\
  &\qquad \text{\olref{L-3} توں شرط P2 نال}\notag \\
  & \OProv[\Th{T}](\gn{!D}) \lif (\OProv[\Th{T}](\gn{\OProv[\Th{T}](\gn{!D})}) \lif \OProv[\Th{T}](\gn{!A})) \ollabel{L-5}\\
  &\qquad \text{\olref{L-4} توں دوبارہ P2 نال} \notag\\
& \OProv[\Th{T}](\gn{!D}) \lif \OProv[\Th{T}](\gn{\OProv[\Th{T}](\gn{!D})}) \ollabel{L-6}\\
  & \qquad\text{!!{derivability} شرط P3 نال} \notag\\
  & \OProv[\Th{T}](\gn{!D}) \lif \OProv[\Th{T}](\gn{!A}) \ollabel{L-7} \\
  &\qquad\text{\olref{L-5} تے \olref{L-6} توں}\notag\\
  & \OProv[\Th{T}](\gn{!A}) \lif !A \ollabel{L-8}\\
  &\qquad\text{قضیے دے مفروضے توں} \notag\\
  & \OProv[\Th{T}](\gn{!D}) \lif !A \ollabel{L-9}\\
  &\qquad\text{\olref{L-7} تے \olref{L-8} توں}\notag\\
  & (\OProv[\Th{T}](\gn{!D}) \lif !A) \lif !D \ollabel{L-10}\\
  & \qquad \text{\olref{L-1} توں}\notag \\
  & !D \ollabel{L-11}\\
  & \qquad\text{\olref{L-9} تے \olref{L-10} توں}\notag \\
  & \OProv[\Th{T}](\gn{!D}) \ollabel{L-12}\\
  & \qquad\text{\olref{L-11} توں شرط~P1 نال}\notag \\
  & !A \qquad\qquad\text{\olref{L-8} تے \olref{L-12} توں}\notag
\end{align}
\end{proof}

لوب دے قضیے توں نامکملیت دے دوجے قضیے دا اک مختصر ثبوت
وی ملدا اے،اوہناں نظریاں لئی جنہاں دا !!a{derivability}
محمول P1--P3 شرطاں پوری کردا ہووے: جے $\Th{T} \Proves
\OProv[\Th{T}](\gn{\lfalse}) \lif \lfalse$،تاں
$\Th{T} \Proves \lfalse$۔ جے $\Th{T}$ بے تضاد اے،
تاں $\Th{T} \Proves/ \lfalse$۔ ایس لئی
$\Th{T}
\Proves/ \OProv[\Th{T}](\gn{\lfalse}) \lif \lfalse$،
یعنی $\Th{T} \Proves/
\OCon[\Th{T}]$۔ اسی قضیے نال ایہہ وی دکھا سکدے آں کہ
$!H$،جیہڑا $\OProv[\Th{T}](x)$ دا مقرر نقطہ اے،
!!{derivable} اے۔ کیونکہ
\begin{align*}
  \Th{T} & \Proves \OProv[\Th{T}](\gn{!H}) \liff !H\\
  \intertext{خاص طور اُتے}
    \Th{T} & \Proves \OProv[\Th{T}](\gn{!H}) \lif !H
\end{align*}
تے ایس لئی لوب دے قضیے نال $\Th{T} \Proves !H$۔

% Going in the other direction, for homework I
% may ask you to work through a short proof of L\"ob's theorem, using
% the second incompleteness theorem instead of the fixed-point lemma.

\begin{prob}
منّو $\Th{T}$ اک قابلِ حساب طریقے نال مسلّمہ بند نظریہ
اے،تے $\OProv[\Th{T}]$، $\Th{T}$ دا
!!a{derivability} محمول اے۔ ہِٹھلے چار بیان ویکھو:
\begin{enumerate}
\item جے $T \Proves !A$،تاں $T \Proves \OProv[\Th{T}](\gn{!A})$۔
\item $T \Proves !A \lif \OProv[\Th{T}](\gn{!A})$۔
\item جے $T \Proves \OProv[\Th{T}](\gn{!A})$،تاں $T \Proves !A$۔
\item $T \Proves \OProv[\Th{T}](\gn{!A}) \lif !A$
\end{enumerate}
ہر اک بیان کنہاں شرطاں وچ درست اے؟
\end{prob}

\end{document}
